\subsection{Dual of a hilbertian space}

THEOREM 3. --- Let E be a hilbertian space. For every \(x \in E\) , let \(x^{*}\) be the continuous linear form \(y \mapsto \langle x | y \rangle\) on E; the mapping \(x \mapsto x^{*}\) is a bijective, semi-linear (for the automorphism \(\xi \mapsto \overline{\xi}\) ) mapping from E onto its dual \(E'\) , and an isometry from the normed space E onto the normed space \(E'\) .

The mapping \(x \mapsto x^*\) is semi-linear by (2) (V, p.~1) and by virtue of the Cauchy-Schwarz inequality, we have \(\| x^* \| = \sup_{\| y \| \leqslant 1} |\langle x | y \rangle| = \| x \|\) , hence \(x \mapsto x^*\) is an isometry from E into E', and in particular, is injective. To complete the proof, we need to prove that for all \(x' \neq 0\) in E', there exists \(x \in E\) such that \(x' = x^*\) . But the hyperplane \(H = Ker x'\) is closed in E; its orthogonal is a line D. Let \(b\) be a non-zero element of D; the kernel of the linear form \(b^*\) is equal to H and hence there exists a scalar \(\lambda \neq 0\) such that \(x' = \lambda . b^* = (\overline{\lambda}. b)^*\) . Q.E.D.

The mapping \(x \mapsto x^{*}\) from E onto its dual \(E'\) is said to be canonical. The inverse mapping from \(E'\) onto E is also called canonical and is denoted by \(x' \mapsto x'^{*}\) . We have

\[
\langle x | y \rangle = \langle y, x ^ {*} \rangle , \quad \langle x, x ^ {\prime} \rangle = \langle x ^ {\prime *} | x \rangle\tag{17}
\]

for \(x, y\) in E and \(x'\) in E'. Also \((x^{*})^{*} = x\) for \(x \in \mathbf{E}\) .

When K is R, the mapping \(x \mapsto x^{*}\) is linear. We shall transfer the scalar product of E to \(E'\) by this mapping. When \(K = C\) , we can consider the mapping \(x \mapsto x^{*}\) as an isomorphism from the vector space \(\overline{E}\) , the conjugate of E onto \(E'\) (V, p.~6). We shall transfer the scalar product of \(\overline{E}\) to \(E'\) by this mapping.

In the two cases considered, \(E'\) is a hilbertian space and we have the formulae

\[
\langle x ^ {*} | y ^ {*} \rangle = \overline {{\langle x | y \rangle}}, \quad \langle x ^ {\prime} | x ^ {\prime} \rangle = \| x ^ {\prime} \| ^ {2}
\]

for \(x, y\) in \(\mathbf{E}\) and \(x'\) in \(\mathbf{E}'\) .

To say that the vector \(x \in E\) is orthogonal to a vector \(y \in E\) is equivalent to saying that the linear form \(x^{*} \in E'\) is orthogonal to y in the sense defined in II, p.~41 (this justifies the use of the word « orthogonal » in the two cases). If M is a closed vector subspace of E, the subspace \(M^{\circ}\) orthogonal to M in \(E'\) (II, p.~44) is the image under \(x \mapsto x^{*}\) of the orthogonal of M in E, defined in V, p.~13 (this justifies the use of the notation \(M^{\circ}\) in the two cases).

COROLLARY 1.--- In order that the family \((x_{i})_{i\in I}\) of points of a hilbertian space E be total, it is necessary and sufficient that the relations \(\langle x_{i}|y\rangle = 0\) for \(y \in E\) and for all indices \(i \in I\) imply that y = 0.

In fact, this says that 0 is the only vector of \(\mathbf{E}'\) which is orthogonal to all the \(x_{i}\) (II, p.~43 and IV, p.~1).

COROLLARY 2. --- Let E and F be two hilbertian spaces. For \(u \in \mathcal{L}(\mathrm{E};\mathrm{F})\) , \(x \in \mathrm{E}\) and \(y \in \mathrm{F}\) , put

\[
\Phi_ {u} (y, x) = \langle y | u (x) \rangle .\tag{18}
\]

The mapping \(u \mapsto \Phi_u\) is an isomorphism from the Banach space \(\mathcal{L}(\mathrm{E};\mathrm{F})\) onto the space of all continuous sesquilinear \(^1\) forms on \(\mathrm{F} \times \mathrm{E}\) , endowed with the norm

\[
\| f\| = \sup_{\substack{x\in \mathbf{E},  y\in \mathbf{F}\\ \| x\| \leqslant 1, \| y\| \leqslant 1}}\left|f(y,x)\right|.\tag{19}
\]

It is clear that \(\Phi_u\) is sesquilinear and continuous for all \(u \in \mathcal{L}(\mathrm{E};\mathrm{F})\) . Conversely, let \(f\) be a continuous sesquilinear form on \(\mathrm{F} \times \mathrm{E}\) . For every \(x \in \mathrm{E}\) , the mapping \(y \mapsto \overline{f(y,x)}\) is a continuous linear form on the hilbertian space \(\mathrm{F}\) . By th. 3, for every \(x \in \mathrm{E}\) , there exists a unique element \(u(x)\) in \(\mathrm{F}\) such that \(\overline{f(y,x)} = \langle u(x)|y\rangle\) for all \(y \in \mathrm{F}\) . The mapping \(u: x \mapsto u(x)\) from \(\mathrm{E}\) into \(\mathrm{F}\) is linear and we have

\[
\begin{array}{l} \| f \| = \sup _ {\| x \| \leqslant 1} \sup _ {\| y \| \leqslant 1} | f (y, x) | = \sup _ {\| x \| \leqslant 1} \sup _ {\| y \| \leqslant 1} | \langle y | u (x) \rangle | \\ = \sup _ {\| x \| \leqslant 1} \| u (x) \|; \end{array}
\]

hence \(u\) belongs to \(\mathcal{L}(\mathrm{E};\mathrm{F})\) , \(f = \Phi_u\) and \(\| u\| = \| f\|\) . This proves cor. 2.

\(^{1}\) Recall (A, IX, § 1, No.~5) that a sesquilinear form (on the left) f on F × E is a mapping from F × E into K which satisfies relations (1) and (2) of V, p.~1.

The canonical mapping from E into its bidual \(E''\) (IV, p.~14) maps E onto \(E''\) , in other words (IV, p.~16), E is a reflexive Banach space. In fact, if E is a real (resp. complex) hilbertian space, the canonical mapping \(\phi\) from \(E'\) onto E is an isomorphism from the normed space \(E'\) onto E (resp. onto the conjugate space \(\overline{E}\) of E); applying th. 3 to E (resp. \(\overline{E}\) ), we see that every continuous linear form on the normed space \(E'\) is of the form \(x' \mapsto \langle \phi(x') | x \rangle = \langle x, x' \rangle\) with \(x \in E\) , hence our assertion follows.

As a consequence (IV, p.~17, prop. 6):

THEOREM 4. --- In a hilbertian space E, the unit ball is weakly compact.

PROPOSITION 10. --- If, in a hilbertian space E, a filter F converges weakly to \(x_{0}\) , and if moreover \(\lim_{\mathfrak{F}} \|x\| = \|x_{0}\|\) , then F converges to \(x_{0}\) for the initial topology of E. In fact, \(\|x - x_{0}\|^{2} = \|x\|^{2} - 2R \langle x|x_{0} \rangle + \|x_{0}\|^{2}\) . Since \(\langle x|x_{0}\rangle\) tends to \(\|x_{0}\|^{2}\) with respect to F by hypothesis, and \(\|x\|\) tends to \(\|x_{0}\|\) with respect to F, \(\|x - x_{0}\|\) tends to 0 with respect to F, hence the proposition.

Remark. --- If E is a Hausdorff prehilbertian space and \(\hat{E}\) the hilbertian space completion of E, we know (III, p.~16) that the dual \(E'\) of E can be identified with the dual of \(\hat{E}\) ; it then follows from th. 3 (V, p.~15) that every continuous linear form on E can be written in a unique way as \(x \mapsto \langle a|x\rangle\) , where \(a \in \hat{E}\) .

